\section{\acs{PoC} Network Application}
\label{chap:proposal:netapp}

The exposure interfaces provided by the testbed through \acs{NEF} and CAMARA northbound \acsp{API}, and protected by \acs{CAPIF}, will support the development and validation of future network applications. With them, developers can interact programmatically with the network and leverage its capabilities to provide services that cater to industry verticals and their associated use cases. In Section \ref{sota:networkapplications:usecases}, there have been highlighted a series of possible network applications across distinct vertical industries and scenarios, such as \acl{CAM}, \acl{T&L}, Video Streaming and \acf{PPDR}.

Among these, \acs{PPDR} emerges as the most demanding in terms of service requirements, as failures in communication or service continuity may compromise the outcome of emergency operations, with potentially irreversible consequences. More specifically, scenarios such as \acs{IOPS}~\cite{IOPS5G} and \acs{FIDEGAD}~\cite{FIDEGAD5G} impose strict requirements on communication continuity, latency, and bandwidth, under dynamic and often degraded network conditions. This makes \acs{PPDR} an ideal vertical context for a \acl{PoC}, where a network application operating under such conditions would necessarily have to be deeply integrated with the \acs{QoS}, Monitoring, and Analytics interfaces provided by the testbed to securely adapt to network changes and enforce the service guarantees that mission-critical services demand.

The \acs{PoC} application is designed within the \acs{FIDEGAD} scenario~\cite{FIDEGAD5G}, in which drone-based fire detection and ground assistance capabilities are deployed to support emergency response teams in wildfire situations. Inherently, the envisioned network application consists of two main components: (i) a camera client, intended to be deployed at the drone (i.e. the \acs{UE}), that captures live video and streams it over a \acs{5G} connection; and (ii) a server, intended to be deployed at the \acs{MEC} or the cloud, that receives the stream, processes it through an \acs{AI} fire detection model and presents the results through a web interface accessible to ground operators, raising an alert whenever fire is identified in a frame. Given the mission-critical nature of this scenario, the network application must sustain continuous video delivery and respond dynamically to changes in network conditions, adjusting its \acs{QoS} profile as handovers are performed between cells, which naturally occurs as the drone surveys different geographical areas in search for fire outbreaks.

\subsection{Application Requirements}
\label{chap:proposal:netapp:requirements}

With the context and purpose of the network application being established, the next step is to derive the requirements for the \acs{PoC} application. These will be organized into the following categories: functional, performance and non-functional requirements. Functional requirements describe the system and its functionality, i.e. the tasks it is capable of performing. Performance requirements define the extent, and under what conditions, such functions must be executed, expressed as quantitative and individually verifiable metrics. Lastly, non-functional requirements specify requirements under which the system is required to operate, as well as the system's properties. Within this category, the focus will be placed on quality requirements, which capture system criteria such as reliability, maintainability, security and portability~\cite{ISORequirementsEngineering}.

\subsubsection{Functional Requirements}
\label{chap:proposal:netapp:requirements:functional}

The functional requirements align with the primary functionalities provided by the \acs{FIDEGAD} network application. These encompass the streaming and processing of the drone's camera feed with the goal of detecting and reporting fire outbreaks. Moreover, it includes the enforcement of \acs{QoS} profiles as the drone (i.e. the \acs{UE}) moves between \acs{5G} cells and their corresponding coverage areas, ensuring the network application maintains the necessary \acs{QoS} to support continuous and uninterrupted operation.

\begin{custenumerate}{FR}
    \item The system shall continuously transmit the drone camera video stream to the processing server over the \acs{5G} network;
    \item The system shall analyze each frame as it is received from the video stream to detect the presence of fire outbreaks;
    \item The system shall notify ground operators upon detection of fire outbreaks in the video stream;
    \item The system shall determine whether the drone (i.e. the \acs{UE}) has transitioned to a new operational area covered by a different \acs{5G} cell, displaying the current serving cell to the ground operator;
    \item The system shall update the active \acs{QoS} profile of the drone's \acs{5G} data session upon transitioning to a new operational area, displaying the current \acs{QoS} parameters to the operator in real time;
    \plan{
        What about these?
        \item The system shall subscribe to UE location change events from the 5G network for the drone's associated UE identifier upon mission initialization;
        \item The system shall apply a fallback QoS profile to the drone's 5G data session upon detection of network degradation.
    }
\end{custenumerate}

\subsubsection{Performance Requirements}
\label{chap:proposal:netapp:requirements:performance}

The performance requirements of the \acs{FIDEGAD} network application are directly tied to the \acs{QoS} parameters configured for the drone's data path. As the application's core functionality is dependent on the streaming of a live video over \acs{5G}, the verifiability of each related performance requirement is contingent on the network being configured with a \acs{QoS} profile that provides the necessary resource guarantees. Therefore, the following pre-condition applies to all performance requirements listed in this section:

\begin{custenumerate}{C}
    \item The \acs{QoS} profile assigned to the drone's \acs{5G} data session shall meet the minimum bitrate required to ensure that each encoded video frame is delivered to the processing server within 33 ms of its capture, corresponding to a frame delivery rate of 30 fps (FR1), as determined during system configuration.
\end{custenumerate}

With this pre-condition established, the performance requirements define measurable aspects of the application's operation, restraining the functional requirements specified in Section \ref{chap:proposal:netapp:requirements:functional} to temporal constraints conditioned by the active \acs{QoS} profile and application processing overhead, as follows:

\begin{custenumerate}{PR}
    \item The system shall maintain both downlink and uplink bitrates within the bounds guaranteed by the active \acs{QoS} profile, each remaining above their respective minimum bitrate and below their respective maximum throughput limit;
    \item The system shall produce a detection result for each received video frame prior to the arrival of the subsequent frame (FR2), as constrained by condition C1;
    \item The system shall notify the ground operator of a fire alert (FR3) within 1 second of the detection event, independently of the active \acs{QoS} profile;
    \item The system shall detect a handover within 1 second of the drone transitioning to a new serving cell (FR4), independently of the active \acs{QoS} profile;
    \item The system shall enforce the \acs{QoS} profile within 1 second of the drone completing the handover to the new \acs{5G} cell (FR5).
    \plan{
        What about these?
        \item The system shall trigger a \acs{QoS} adaptation when measured session performance degrades below the minimum thresholds specified by the active \acs{QoS} profile.
        \item Also where are jitter and RTT jitter (we should be using this)
    }
\end{custenumerate}

\subsubsection{Non-Functional Requirements}
\label{chap:proposal:netapp:requirements:nonfunctional}

The following list of non-functional requirements addresses quality attributes of the system that are independent of the specific operational scenarios highlighted earlier. They cover a range of concerns, including:

\begin{custenumerate}{NR}
    \item All northbound \acs{API} interactions shall be authenticated and authorized through a single secure exposure layer;
    \item The system shall recover autonomously from transient disconnections between the client and server, resuming video transmission without operator intervention;
    \item The system shall produce structured logs of all \acs{API} interactions, fire detection events and \acs{QoS} profile changes to support debugging and validation procedures;
    \item Both application components shall be packaged as \acs{NFV} or containerized applications, deployable through the testbed's orchestration infrastructure.
\end{custenumerate}

\subsection{Interface Requirements}

\plan{
    \begin{enumerate}
        \item Introduce the PPDR domain as the vertical context for the PoC, drawing from the use cases discussed in Section \ref{sota:networkapplications:usecases} (IOPS, FIDEGAD). Highlight the particular relevance of PPDR scenarios for exercising network exposure interfaces: they require reliable, low-latency communication under dynamic and potentially degraded network conditions, making QoS enforcement and continuous network monitoring essential, not optional, capabilities.
        \item State the objective of the \acs{PoC} application: a Public Protection and Disaster Relief (PPDR) network application developed within the FIDEGAD scenario, where drone-based fire detection and ground assistance capabilities are leveraged to support emergency response teams in wildfire situations. Reenforce this as an ongoing mission-critical operation, dynamically adapting QoS profile in response to changes in network conditions or operational context.
        \item Derive the application requirements from the use case and divide them between Functional Requirements and non functional requirements. Here we have some examples of what's expected:
        \begin{itemize}
            \item Guaranteed minimum UL/DL throughput for video data transmission;
            \item Bounded RTT and low jitter for real-time command and control;
            \item Continuous monitoring of network performance against these requirements;
            \item Ability to escalate to a higher-priority QoS profile when degradation is detected or predicted.
        \end{itemize}
        \item The implemented exposure interfaces, provided by the CAMARA and NEF northbound \acsp{API}, and secured by \acs{CAPIF}, will serve as the foundation for the development and validation of future network applications, by a variety of developers.
        \item Revisit the use cases from Section \ref{sota:networkapplications:usecases} and corresponding internal details (\acs{NEF}/CAMARA \acs{API} suite) from Section \ref{sec:testingandvalidation:netapps}:
        \begin{itemize}
            \item Dynamic \acs{QoS} adaptation enforced through CAMARA \acs{QoS} \acs{API} in response to either geographical changes obtained with CAMARA Device Location \acs{API} in digital twin / \acs{CAM} scenarios, or to network quality deterioration (or to video quality deterioration in case of video streaming/broadcast).
            \item \acs{PPDR}: \acs{IOPS}, \acs{FIDEGAD}
        \end{itemize}
        \item Map each requirement to the corresponding network interface, showing how the PoC will interact with the implemented components to fulfill its goals:
        \begin{itemize}
            \item CAPIF: onboarding and authenticated discovery of all APIs;
            \item NEF AsSessionWithQoS API: direct QoS profile enforcement per UE flow;
            \item NEF MonitoringEvent API: subscription to UE state events (e.g. location changes) that may warrant a QoS adaptation;
            \item NEF Analytics Exposure API: retrieval of live throughput, RTT, jitter, and packet loss metrics;
            \item CAMARA Application Profiles API: declaration of the application's network requirements as a machine-readable profile;
            \item CAMARA Connectivity Insights Subscriptions API: event-driven notification when network performance falls below declared thresholds.
        \end{itemize}
        \item Clarify that NEF and CAMARA APIs are not redundant but complementary alternatives: the CAMARA APIs provide a higher-level, developer-friendly abstraction over the same underlying NEF capabilities, while the NEF APIs offer finer-grained control. Both paths are exercised in the PoC to validate the full integration stack.
        \item Include an interaction sequence diagram illustrating the end-to-end flow of the PoC (prioritizing CAMARA):
        \begin{enumerate}
            \item (in progress)
            \item Onboarding via CAPIF
            \item the creation of application profile according declaring application requirements (does not need to match exactly the slice requirements)
            \item the creation of connectivity insights subscription to be notified when the application requirements are not being respected
            \item  the change of \acs{QoS} profile in response to being unable to meet the application requirements or in response to a location change (CAMARA QoD API)
        \end{enumerate}
        \item Explain how the \acs{PoC} exercises all implemented interfaces end-to-end: it onboards via CAPIF, consumes CAMARA APIs (Application Profiles to declare its requirements, Connectivity Insights Subscriptions to monitor performance, QoD to enforce them)
        \item Include another sequence diagram this time prioritizing the NEF. Showcase the added complexity of the NEF implementation flow comparatively to CAMARA APIs.
    \end{enumerate}
}
